% !TEX TS-program = lualatex

\DocumentMetadata{
  lang = en,
  pdfversion = 2.0,
  pdfstandard = ua-2,
  tagging = on,
  tagging-setup = {
    math/setup = {mathml-SE},
    table/header-rows = 1
  }
}

\documentclass[12pt]{article}

\usepackage{amsmath}
\usepackage{geometry}
\geometry{
  letterpaper,
  textwidth=6.5in,
  textheight=9in,
  top=0.85in,
  bottom=1in,
  left=0.75in,
  right=0.75in,
}
\usepackage{hyperref}

\newtheorem{theorem}{Theorem}
\newtheorem{acknowledgement}[theorem]{Acknowledgement}
\newtheorem{algorithm}[theorem]{Algorithm}
\newtheorem{axiom}[theorem]{Axiom}
\newtheorem{case}[theorem]{Case}
\newtheorem{claim}[theorem]{Claim}
\newtheorem{conclusion}[theorem]{Conclusion}
\newtheorem{condition}[theorem]{Condition}
\newtheorem{conjecture}[theorem]{Conjecture}
\newtheorem{corollary}[theorem]{Corollary}
\newtheorem{criterion}[theorem]{Criterion}
\newtheorem{definition}[theorem]{Definition}
\newtheorem{example}[theorem]{Example}
\newtheorem{exercise}[theorem]{Exercise}
\newtheorem{lemma}[theorem]{Lemma}
\newtheorem{notation}[theorem]{Notation}
\newtheorem{problem}[theorem]{Problem}
\newtheorem{proposition}[theorem]{Proposition}
\newtheorem{remark}[theorem]{Remark}
\newtheorem{solution}[theorem]{Solution}
\newtheorem{summary}[theorem]{Summary}
\newenvironment{proof}[1][Proof]{\textbf{#1.} }{\ \rule{0.5em}{0.5em}}

\tagpdfsetup{
  math/alt/use,
  role/new-tag=section/H1,
}

\hypersetup{
  pdftitle={Assignment 6},
  pdfauthor={Blankenau},
  pdfsubject={Fall, 2025},
  hidelinks,
}

\begin{document}
\begin{center}
{\LARGE Assignment 6}
\end{center}

\noindent Blankenau

\noindent Economics 905, Fall, 2025

\medskip
\noindent \noindent \textit{Instructions.} You may work in groups and may
compare answers prior to submission. Joint submissions with up to 4 names
are allowed and encouraged. All work should be legible and concise. A due
date will be announced in class.

\begin{enumerate}
\item Describe how each of the following affects the $\dot{k}=0$ and $\dot{c}%
=0$ curves in our phase diagram, and thus how they affect the
balanced-growth-path values of $c$ and $k$:

\begin{enumerate}
\item A rise in $\theta .$

\item A change in the depreciation rate from the value of 0 assumed so far
to some positive level. That is, capital now evolves according to $\overset{.%
}{k}\left( t\right) =f\left( k\left( t\right) \right) -c\left( t\right)
-\left( n+g+\delta \right) k\left( t\right) .$ Note in addition that now the
return on capital is $r\left( t\right) =f^{\prime }\left( k\left( t\right)
\right) -\delta \ $rather than $r\left( t\right) =f^{\prime }\left( k\left(
t\right) \right) .$ This should lead you to conclude that both the $\dot{k}%
=0 $ and $\dot{c}=0$ shift.
\end{enumerate}

\item Consider a standard Ramsey-Cass-Koopman's growth model where
investment income is subsidized at rate $\varsigma _{k}.$ In each period,
revenue for the subsidy is generated from a lump sum tax of $G\left(
t\right) $ per effective unit of labor. The representative household
maximizes 
\begin{equation*}
U=\dint\limits_{t=0}^{\infty }e^{-\left( \rho -n-\left( 1-\theta \right)
g\right) t}\frac{c\left( t\right) ^{1-\theta }}{1-\theta }dt
\end{equation*}%
$\theta >0,$ $\ \rho -n-\left( 1-\theta \right) g>0$ subject to%
\begin{equation*}
k\left( 0\right) +\dint\limits_{t=0}^{\infty }e^{-R\left( t\right)
}e^{\left( n+g\right) t}\left( w\left( t\right) -G\left( t\right) \right)
dt=\dint\limits_{t=0}^{\infty }e^{-R\left( t\right) }e^{\left( n+g\right)
t}c\left( t\right) dt
\end{equation*}%
and the standard transversality condition. The capital stock per effective
unit of labor evolves according to $\overset{.}{k}\left( t\right) =f\left(
k\left( t\right) \right) -c\left( t\right) -\left( n+g\right) k\left(
t\right) .$ All items are defined as in the text: $c\left( t\right) $ and $%
k\left( t\right) $ are consumption and capital per effective unit of labor, $%
n$ in the growth rate of the population, $g$ is the growth rate of the labor
augmenting technology, $\rho $ is the discount rate, $R\left( t\right)
=\dint\limits_{\tau =0}^{t}r\left( \tau \right) d\tau $ $\ $where $r\left(
\tau \right) $ is the period $\tau $ interest rate, $w\left( t\right) $ is
the period $t$ wage. Thus the real interest rate facing households is given
by $\left( 1+\varsigma _{k}\right) f^{\prime }\left( k\left( t\right)
\right) .$ The intensive form production function $f\left( k\left( t\right)
\right) $ has the standard properties.

\begin{enumerate}
\item Show your derivation of the optimal growth rate of consumption.

\item Draw the $\overset{.}{k}\left( t\right) =0$ and $\overset{.}{c}\left(
t\right) =0$ loci in the $c,\,k$ space (that is, draw the principal curves
of the phase diagram). Label everything properly. Label the steady state
level of capital $k^{\ast }.$

\item For this part only let $n=g=0$ and $f\left( k\left( t\right) \right)
=y\left( t\right) =\left( k\left( t\right) \right) ^{.5}$. Suppose that the
economy has been operating at $k\left( t\right) =\frac{K_{0}}{L_{0}A_{0}}%
=k_{0}.$ Suddenly $k\left( t\right) $ changes to $k=zk_{0}\ $where $z<1$
This might have happened because of destruction of some of the capital
stock, a one-time permanent increase in $L$ or a one-time permanent increase
in $A$. A policy maker says that it does not matter what caused this
decrease since all dynamics are the same, regardless of the cause of the
decrease. Do you agree with this? If so, explain why. If not, describe the
differences across the three cases, paying particular attention to $y\left(
t\right) $ and per capita output. Consider these items both immediately and
in the long run.

\item Suppose the economy is again at $k=k^{\ast }.$ What are the immediate
and long run consequence on $c$ and $k$ of an increase in $\varsigma _{k}?$
Show this on a phase diagram.

\item A policy maker says that in this situation, the effect of an increase
in $\varsigma _{k}=0$ to $\varsigma _{k}>0$ on utility, $U$, is ambiguous
because consumption is affected differently in the short run and in the long
run. Do you agree with this statement? Why or why not?
\end{enumerate}
\end{enumerate}
\end{document}
